% tex/sections/results.tex
% ──────────────────────────────────────────────────────────────────────────────
\chapter{Results}
\label{chap:results}
% ──────────────────────────────────────────────────────────────────────────────
%
% TODO: Write results.
%
% Suggested structure:
%   - Present findings in logical order (mirrors methodology)
%   - Figures and tables with captions
%   - Keep interpretation minimal here — save it for Discussion

\section{Simple Model Results}
\label{sec:results-1}

The following results shown are for the supershear and sub-rayleigh case of the custom stfs. The sources are spaced 0.1m apart between 30m and 270m and have a ricker wavelet all directions. The moment tensors have components 1 in xx and xy and -1 in yy. The simple, two layer model was set up to test the custom STF option in SALVUS and to verify how it works. The results are shown here because it is important to understand the physics underlying crack propagation in snow before moving on to the more complex case. Since the focus of this thesis is to model crack propagation as recorded by DAS data, strain and strain rate plots are shown first. The velocity results are shown to verify observations made in the strain and strain rate plots. The results of the CZM model will not be discussed in detail. This is because the simple model results suffice to understand the MPM-aided model results. Additionally, if physics of the CZM model are kept simple, for example when there is no mode mix of the fracturing, the CZM results are very close to the simple model results. Nevertheless, setting up this model was important to understand snow mechanics and crack propagation physics. \par

\subsection{Sub-rayleigh strain and strain rate plot}
\begin{figure}[H]
    \centering
   \includegraphics[width=\textwidth]{preliminary_sub_strain.png}
   \caption{Sub-rayleigh strain and strain rate (horizontal component) averaged over 10.4 m gaugue length to represent DAS data: strain rate in top column, strain in bottom column. Crack depth on left side, snow bottom on right side. Colorbar shows magnitude of strain or strain rate, colorbars unscaled.}
   \label{fig:sub_strain}
\end{figure}

Both the strain and the strain rate at crack depth look very similar to each other, the same is true for the snow bottom. At crack depth, both the strain and strain rate in the left column of figure \ref{fig:sub_strain} show one strong main front, which is preceded by a preliminary front. The main front with higer amplitude has a slope which corresponds to the prescribed rupture speed, whereas the precursor is faster and has lower amplitude (fainter color). The same is true for the snow bottom, but here, the main front is thinner than the preliminary front. The main front still has higher amplitude than the preliminary front, but the preliminary front weakens in amplitude with increasing time. \par

At crack depth, the main front broadens with increasing time, this can not be observed for the main front measured at the snow bottom on the right of figure \ref{fig:sub_strain}. At both depths, for strain and strain rate, the two fronts have visibly different slopes, which corresponds to different apparent propagation speeds. \par

At both depths, a continous wavefront extends from 0s starting at 30m to 6s ending at 270m, this corresponds to the source coordinates defined and also to the time delays prescribed to the sources. \par

In all plots, the strains and strain rates following each other have opposite polarity. The magnitude of strain and strain rate is larger at the crack depth compared to the bottom of the snow. This is indicated by the colourbar range at the snow bottom being approximately one order of magnitude smaller than at crack depth. This indicates strong amplitude decay over the 0.375,m vertical separation.\par

At x=270m, t=6s in all sub-rayleigh plots in figure \ref{fig:sub_strain}, the main wavefront stops, but there is a new wavefront propagating from this point. This front has a change in slope, which means a different propagation speed, which coincides with the rupture arrest front. \par 


\subsection{Supershear strain and strain rate plot}
\begin{figure}[H]
    \centering
   \includegraphics[width=\textwidth]{preliminary_super_strain.png}
   \caption{Supershear strain and strain rate (horizontal component) averaged over 10.4 m gaugue length to represent DAS data: strain rate in top column, strain in bottom column. Crack depth on left side, snow bottom on right side. Colorbar shows magnitude of strain or strain rate, colorbars unscaled.}
   \label{fig:super_strain}
\end{figure}

In the supershear case, strain and strain rate at both receiver depths in figure \ref{fig:super_strain} show a single dominant wavefront with comparable amplitude, with no significant reduction between crack depth and snow bottom. This main wavefront has strong amplitude at first, but experiences significant reduction in amplitude after a while. This corresponds to loss of energy of the different frequencies contained in the wavefront. \par

The strain and strain rates observed have opposite polarity (positive followed by negative). In the strain rate plots, the different polarities can be seen a bit more clearly than in the strain plots in the bottom row of figure \ref{fig:super_strain}. \par

In figure \ref{fig:super_strain}, there is no difference in the magnitude of strain and strain rate between the crack depth and snow bottom. However, the peak strain rate in the supershear case is approximately one order of magnitude larger than in the sub-Rayleigh case, as indicated by the colourbar scales in figures \ref{fig:sub_strain} and\ref{fig:super_strain}. \par

Similarly to the sub-rayleigh plots in figure \ref{fig:sub_strain}, the supershear plots also show wavefronts with different amplitudes at the rupture arrest at 270m. Here, this is characterized as a fan of distinct wave trains coming from the rupture arrest point. \par


\subsection{Sub-Rayleigh Velocity plots}
The following plots show horizontal (vx) and vertical (vy) particle velocity at two reciever depths: 2.25m, which is the snow bottom, and 2.625m which is the depth of the crack. The particle velocity shown is derived from the displacement field, this is done to assimilate the results shown more to DAS field measurements. In DAS field data, the instantaneous particle velocity at one point is not measured, but rather avegraged over a gauge length \parencite{lindseyFiberOpticSeismology2021}. By taking the temporal derivative of displacement, the results shown here are closer to field results. \par 

\begin{figure}[H]
    \centering
   \includegraphics[width=\textwidth]{preliminary_sub_vel.png}
   \caption{Sub-Rayleigh Particle velocity: vertical velocity in top row, horizontal in bottom row. Crack depth in left, snow bottom in right column. Colorbar is unscaled, showing velocity (m/s).}
   \label{fig:sub_vel}
\end{figure}

The sub-rayleigh velocity plots in figure \ref{fig:sub_vel} all show a distinct wavefront starting at 0s at 30m and ending at 270m. This main wave train ends at 6s, but sends out waves with a different speed (different angle) in direction of the source propagation. This corresponds to the arrest of the rupture which was already observed in the strain and strain rate plots. The horizontal velocity at the snow bottom also shows a backward-propagating train from the crack tip. \par

Both vx plots (bottom row in figure \ref{fig:sub_vel}), show the main wavefront with waves travelling ahead of it, as time and distance progresses, the main wavefront gets wider. This means that they have different slopes, which corresponds to different propagation velocities. In all but the plot of vx at the snow bottom, the first couple of waves in the main wavefront are weaker in color, and thus weaker in velocity. In vx at the crack depth at the bottom left of figure \ref{fig:sub_vel}, the precursor is weaker in color and thus magnitude of velocity than the main wavefront. At the snow bottom, which is shown on the bottom right, the precursor has a high amplitude of velocity initially, which then decays. \par

In the main wavefront observed, it's very noticeable that the particle velocity is layered. In other words, there are stripes along the wavefront with opposite polarity, for example, there is a stripe with positive polarity (red) followed by one of negative ploarity (blue). This is also noticeable in the cases that have multiple wavefronts, although it's less pronounced. \par

In all panels of figure \ref{fig:sub_vel}, the main wavefront between 30m and 270m does not decay in amplitude with increasing time. It is noticeable however that in all but the horizontal velocity at the snow bottom, the main wavefront broadens with time, which corresponds to dispersive wave propagation in a bounded medium. \par



\subsection{Supershear Velocity Plots}
\begin{figure}[H]
    \centering
   \includegraphics[width=\textwidth]{preliminary_super_vel.png}
   \caption{Supershear Particle velocity: vertical velocity in top row, horizontal in bottom row. Crack depth in left, snow bottom in right column. Colorbar is unscaled, showing velocity (m/s).}
   \label{fig:super_vel}
\end{figure}

In all of the supershear velocity plots in figure \ref{fig:super_vel}, a single dominant wavefront extends from 30m to the end of the medium at 400m. The slope of this front corresponds to the prescribed supershear rupture velocity. This wavefront broadens with increasing propagation time, where it also loses amplitude (weaker color). At the crack arrest at 270m, the end of the source array, the wavefront generates a fan of outward-radiating waves consistent with stopping-phase radiation from the crack arrest point \parencite{madariagaHighfrequencyRadiationCrack1977}. \par

In the vy panel at the snow bottom (top right in figure \ref{fig:super_vel}), there is a clear distinction between two wavefronts, which cannot be seen in all other supershear velocity space-time plots. The first wavefront with high amplitude corresponds in slope and therefore propagation speed to the wavefront in all other plots, it also has a fanning out of waves from 270m, so the rupture arrest. From the crack onset at 30m, there is a second wavefront which also is of a higher amplitude. This wavefront has a shallower slope and therefore a lower propagation speed than the main wavefront and does not emit a fan at the rupture arrest point. \par

The wavefronts arriving in all cases switch between polarities. While the pattern of polarity is the same in all cases, the slope of the wavefronts change. In all plots, after the first wavefront, the waves seem to slow, which means their slope flattens. \par

\subsection{Frequency-Wavenumber (fk) Plot}
\label{sec:fk_results}

The fk-plot in figure \ref{fig:preliminary_fk} shows the energy (colorbar) density at frequencies (y-axis) and wavenumbers (x-axis) present in the results. This energy was taken from the horizontal velocity results at the snow bottom because it corresponds to the data measured by a horizontal DAS cable buried in snow. The plots shown are normalized to the 99th percentile of the energy present, this is done so that outliers with large energies are scaled and do not skew the color scale. 


\begin{figure}[H]
    \centering
    \includegraphics[width=\linewidth]{../figures/preliminary_fk.png}
    \caption{F-k plot for the supershear (left) and sub-rayleigh (right) case. Slab p-wave velocity in red, s-wave velocity in turquoise and respective rupture speed in blue. Colorbar shows energy normalized to the 99th percentile.}
    \label{fig:preliminary_fk}
\end{figure}

The supershear and sub-Rayleigh f-k plots in figure \ref{fig:preliminary_fk} show qualitatively different energy distributions, consistent with the two distinct propagation regimes. The supershear case shows a concentration of energy around the rupture speed (blue line) and the vp (red line) of the slab on both sides of the wavenumber spectrum. It is important to note here that these speeds are very close in value, so the lines cannot be distinguised well in the plots. The sub-rayleigh case also shows two concentrations of energy. A very broad band at the vp of the slab and a concentration of energy at the rupture speed on the postive wavenumber side. In all cases, the energy is more concentrated around the vp and or rupture speed for lower frequencies, at higher frequencies, the concentration of energy starts to deviate from speeds present in the medium in figure \ref{fig:preliminary_fk}. \par

In the supershear case, the dark spots of concentrated energy are more broad on the positive side compared to the negative side, where they are narrower and less smeared out. This asymmetry reflects the directionality of the source: the crack propagates in positive x and more energy is released in that direction. In the sub-rayleigh case on the right side of figure \ref{fig:preliminary_fk}, there is no concentration of energy at negative wavenumbers. \par


\section{MPM Guided Results}
\label{sec:mpm_results}

The results shown here are from the MPM-guided simulation. The sources are spaced 0.01m apart from 75m to 125m. Unlike the results in section \ref{sec:results-1}, which are for a simple, 2 layer model with prescribed rupture velocities (sub-rayleigh and supershear case), the results here are for a 5-layer model with a rupture velocity coming from the MPM simulations. The results shown here are filtered using a low-pass filter with 50Hz cutoff. Because both strain and velocity are derived from displacement, the displacement field extracted is shown in appendix \ref{app:I}. The displacement data itself is unfiltered, but because strain, strain rate and velocity are derivatives, the differentiation acts as a high-pass filter which amplifies noise. Displacement is shown in the appendix to show that there are little numerical artefacts and noise in the extracted data. This means that filtering strain, strain rate and velocity should not be a concern because the noise is simply from the derivative. 

\subsection{MPM Strain and Strain rate}
\label{sec:mpm_strain}


\begin{figure}[h]
    \centering
    \includegraphics[width=\linewidth]{../figures/mpm_strain.png}
    \caption{Strain and strain rate (horizontal component) averaged over 10.4m gauge length to represent DAS data: strain rate in top coulmn, strain in bottom column. Crack depth on left, snow bottom on the right. Colorbar shows magnitude of strain or strain rate.}
    \label{fig:mpm_strain}
\end{figure}

The strain as well as the strain rate plots in figure \ref{fig:mpm_strain} all show a distinct v-shape initiating from the crack onset point at x=75m. The left arm of the V corresponds to backward-propagating body waves radiated from the crack initiation point, while the right arm corresponds to the propagating crack traveling from the initiation point onwards. \par

Both strain panels (bottom column) show a front of positive strain (red, on top) arriving first, which is followed by a negative strain front. The extension (positive strain) followed by compression could indicate P-waves or the S0 Lamb wave mode. At both the crack depth as well as the snow bottom, in positive x-direction, the red front widens until about x=125m, when there is a jump and the front has about same width as the blue front. This location coincides with the crack arrest location. In the negative x-direction, both the positive and negative front have the same width from the initiation point on to the domain boundary. \par 

The same positive and negative fronts from the crack onset point can be observed in the strain rate plots in the top row of figure \ref{fig:mpm_strain}. Contrary to the strain panels, where the fronts are solid colors, in the strain rate panels, these are much more heterogenous. The forward propagating positive front is much narrower in strain rate than in strain, the same is true for the backward propagating negative front. The wedge-shape to the crack arrest point observed in strain can still be seen in strain rate however. In the strain rate plots there also more than just two wavefronts in the v-shape coming from the crack onset. The main positive and negative fronts are still visible, but there is another positve front under the blue front in the positive directon. In the negative direction, the same is true: under the blue wavefront there is multiple positive and negative layered wavefronts. \par

In the strain plot at crack depth as well as in the strain rate at the same depth, there are more wavefronts additional to the v-shape from the crack onset. Those fronts fan out from the crack initiation point at 75m as well, but they have a much steeper slope than the main front. In contrast to the main front, they are weaker in color meaning amplitdue. These fronts are layered in polarity, meaning a front of positive polarity (red) is followed by one of negative polarity. \par

These weaker amplitude fronts broaden with increasing time, which is consistent with dispersive wave propagation. Their amplitude does not decrease however. In the strain rate at crack depth (y=0.525m, weak layer, shown on top right of figure \ref{fig:mpm_strain}), the colors of the second wavefront are weaker than in the strain panel at the same depth. This means that they have lower amplitude relative to the main wavefront in the strain rate panel. \par

The main wavefront observed in figure \ref{fig:mpm_strain} in all panels does not broaden over time, it also shows no fading in color, which would indicate a decrease in amplitude. \par 

In all panels, there are two red-blue wavefronts coming from the domain side boundaries. These are especially noticeable in the snow bottom strain and strain rate plot (right side of figure \ref{fig:mpm_strain}), because the main front is the only wavefront present there. The wavefronts coming from the side boundaries of the domain have the same slope as the main wavefronts (wavetrain coming from the right boundary has the same slope as backward-propagating front from crack), but they are much weaker in amplitude. \par 

When comparing the strain rate plots (top column) to the strain plots, allthough the main structures are the same, the colorbar scale in the strain rate plots is one order of magnitude (one power of 10) larger than that of the strain plots. This does not depend on depth (snow bottom has the same amplitude as crack depth) but on the type of field shown. \par



\subsection{MPM Velocity}
\label{sec:mpm_vel}

\begin{figure}[h]
    \centering
    \includegraphics[width=\linewidth]{../figures/mpm_vel.png}
    \caption{Particle velocity: vertical velocity in top row, horizontal velocity (vx) in bottom row. Crack depth on left, snow bottom on the right. Colorbar shows magnitude of particle velocity (m/s).}
    \label{fig:mpm_vel}
\end{figure}

The velocity plots in figure \ref{fig:mpm_vel} show that both cases look pretty similar to each other independent of depth. This means that the vertical velocity in the top row, looks almost identcial at crack depth compared to the snow bottom. The vertical velocity plots can be distinguised from the horizontal velocity plots however. \par 

The horizontal velocity in the bottom two panels of figure \ref{fig:mpm_vel} both show a flat v-shape from the crack origin at 75m, similar to that observed in the strain and strain rate plots in figuer \ref{sec:mpm_strain}. The forward propagating front is characterized by a wavetrain of negative velocity (blue), which is the first arrival, followed by a train of positive energy (red) which arrives shortly after. In the direction backwards of propagation, this order is flipped: a positive front arrives followed by a negative front. In the horizontal velocity plots in figure \ref{fig:mpm_vel}, the vertex of the fronts with opposite polarity seems to be shifted. The first front, which is blue in positive direction of propagation and red in negative direction, is at x=75m, which is the crack onset. The vertex of the second v-shape, which has opposite polarity (red in positive, blue in negative direction), is not right below the vertex of the first v-shape. It is shifted in positive direction of propagation, which indicates that it could come from the crack arrest point a x=125m. \par 

From the same vertex, in the vx plots at crack depth (bottom left in \ref{fig:mpm_vel}), fronts with weaker amplitude originate. These wavefronts fan out from around x=125m and they are layered in polarity (positive followed by negative). They also widen as time incerases, which could indicate dispersion. Such secondary wavefronts are not present in the vx plot at the snow bottom. \par 

In the vertical velocity plots in the top row of figure \ref{fig:mpm_vel}, the secondary wavefronts are very strong and can be seen clearly at both depths. They originate in a v-shape from the crack initiation point at x=75m. These wavefronts are very strong in amplitude (saturated colors), and are layered in polarity. The wvaefronts in both the positive and negative direction of propagation widen as time increases and their amplitude decreases, which is indicative of gemetric spreading and dispersion. When looking at the later times in the top panels of figure \ref{fig:mpm_vel}, the blue and red velocities seem to curve towards the origin of the plot (towards x=0). This is another indicator for dispersion of higher frequencies. \par

Allthough the wavefronts from x=75m are most noticeable, there are v-shaped wavefronts which also originate at x=125m which is the crack arrest point. These are mainly noticeable at earlier times, for example between 1s and 2s, because at later times, there are multiple other wavefronts interfering with them. The second v of waves from the crack arrest starts later in time than the one from the crack initiation point, which is intuitive because the crack onset and arrest are not at the same time. The slope of the individual components from the arrest wavefront is similar to the one of the onset. \par

From the crack initiation point at x=75m, in the vertical velocity panels at the top of figure \ref{fig:mpm_vel}, there is a second wavefront which also propagates in the positive and the negative direction. The wavetrains in this front are much thinner than in the front that was described previously, they also have a flatter slope, which means lower speed. The slope of this wavefront is comparable to the one in the vx plots below them. This wavefront also has layered polarity, but here the bands of one polarity are much thinner than in the vx plots or the other wavefronts in the vy plots. \par

In the area where the crack propagates, so between 75m and 125m, at the early times in vy in the top panels of figure \ref{fig:mpm_vel}, the slope of the thinner layer wavefront is different than in the rest of the wavefront. This area has a slope which is steeper than the rest of the wavefront, it is especially noticeable in vy at the snow bottom, where a red line of positive velocity seems to connect the crack onset and arrest point. The slope of this line is likely the crack propagation speed. \par

In the vertical velocity plot at crack depth, one very faint wave from the crack onset point can also be seen. This wavetrain has a different slope than all other wavefronts in the plot, but it is very weak in amplitude (color). \par 


\subsection{MPM f-k plot}

The data used for the f-k plot is the horizontal velocity data, for the reasons described in \ref{sec:results-1}. 

\begin{figure}[h]
    \centering
    \includegraphics[width = \linewidth]{../figures/mpm_fk.png}
    \caption{F-k plot for the MPM-guided simulation. Slab p-wave velocity in red, slab s-wave velocity in turquoise. Weak layer p-wave velocity shown by the white layer, s-wave velocity of WL shown by yellow line. Colorbar shows energy normalized to the 99th percentile to avoid outliers.}
    \label{fig:mpm_fk}
\end{figure}

The frequency-wavenumber plot in figure \ref{fig:mpm_fk} shows a symmetric distribution of energy around the origin. A lot of energy is concentrated around the positive and negative slab and weak layer p-wave velocity, which are indicated by the red (slab vp) and white (wl vp) lines. \par

The concentration of energy is not strictly at those two speeds, the dark spot of concentrated energy is broader and more spread out and therefore not contained by the two lines. \par

In the f-k plot in figure \ref{fig:mpm_fk}, there is no energy concentrated at the s-wave speeds of the slab, which is indicated by the turqoise line, or the vs of the weak layer, shown by the yellow line. \par

Additional to the energy concentrated around the p-wave speeds, there are bands of concentrated energy on the positive and negative side of the wavenumbers. These energy bands are stronger in the positive direction and they have a shallower slope than the s-wave speeds present in the medium. \par 




% Example figure placeholder:
% \begin{figure}[htbp]
%     \centering
%     \includegraphics[width=0.8\linewidth]{placeholder}
%     \caption{Caption text.}
%     \label{fig:placeholder}
% \end{figure}
